% !TEX root = ../tecnologia-en-marcha.tex
% !TEX encoding = UTF-8 Unicode
% =============================================================================
% Título, autores, resumen y palabras clave.
% Texto: primer borrador del artículo (articulo_tecnologia_en_marcha.md), con los
% acentos restaurados.
% =============================================================================

% El título agrega lo que identifica a la investigación (PIA02): la ENIA y las
% personas adultas mayores. La traducción al inglés es una propuesta: la
% plantilla pide evitar traductores automáticos.
\newcommand{\tituloES}{Marco operativo de cumplimiento de la Estrategia Nacional de IA de Costa Rica para el monitoreo domiciliario de personas adultas mayores con sensores PIR}
\newcommand{\tituloEN}{Operational compliance framework for Costa Rica's National AI Strategy in home monitoring of older adults with PIR sensors}

% Una persona autora por bloque: nombre en negrita; afiliación y país; correo;
% ORCID. La plantilla deja lugar para tres.
\newcommand{\autor}[4]{%
  \textbf{#1}\\
  #2\\
  #3\\
  ORCID: #4\par\medskip}

% Los correos, las cédulas y los ORCID no se versionan aquí: el repositorio es
% público. Viven en la versión de Word, que es la que circula entre las personas
% autoras.
\newcommand{\bloqueAutores}{%
  \autor{Laura María Segura Miranda}
        {Universidad CENFOTEC. Costa Rica}
        {\pendiente{correo}}
        {\pendiente{ORCID}}
  \autor{Miguel Ángel Méndez Arias}
        {Universidad CENFOTEC. Costa Rica}
        {\pendiente{correo}}
        {\pendiente{ORCID}}
  \autor{David Gerardo Alfaro Viquez}
        {Universidad CENFOTEC. Costa Rica}
        {\pendiente{correo}}
        {\pendiente{ORCID}}}

\newcommand{\resumenES}{%
La Estrategia Nacional de Inteligencia Artificial de Costa Rica (ENIA 2024-2027)
establece principios de transparencia y explicabilidad, equidad y no
discriminación, y responsabilidad, pero no especifica mecanismos técnicos para
demostrar su cumplimiento. Este artículo presenta un marco operativo que traduce
esos tres principios en nueve requerimientos verificables, pruebas computacionales
y artefactos auditables. Siguiendo la ciencia del diseño, los requerimientos se
derivaron de la ENIA, el Reglamento Europeo de IA y el NIST AI RMF, y se
implementaron tres módulos: explicabilidad mediante SHAP; equidad mediante
diferencias de tasas de verdaderos y falsos positivos entre subgrupos, con umbrales
declarados antes de medir; y trazabilidad mediante una bitácora que registra cada
inferencia y la vincula con su explicación. El marco se aplicó a un clasificador de
actividades basado en sensores de infrarrojo pasivo (PIR), sobre dos conjuntos de
datos distintos del proyecto CASAS: Aruba, con una vivienda, y una serie de nueve
viviendas de personas adultas mayores. En ambos pilotos produjo un expediente de
evidencia completo, con los nueve requerimientos cubiertos, y el sistema evaluado
obtuvo un veredicto de no cumplimiento. Las disparidades aparecieron entre franjas
horarias, tipos de día y, sobre todo, entre viviendas, con 14 de las 29
combinaciones que superaron el umbral en el piloto multi-hogar. Con diez semillas,
el veredicto se mantuvo y 32 de las 33 combinaciones sobre el umbral persistieron
en todas. Los principios seleccionados pueden operacionalizarse mediante pruebas
reproducibles, aunque la escasez de metadatos demográficos en los datos públicos
limita la evaluación de equidad.}

\newcommand{\palabrasclaveES}{%
Inteligencia artificial; gobernanza algorítmica; cumplimiento normativo de IA;
equidad algorítmica; explicabilidad; sensores de infrarrojo pasivo; envejecimiento
en el lugar (aging in place)
\revisar{la plantilla pide términos del tesauro de la UNESCO; esta es la lista que propuso el profesor}}

\newcommand{\resumenEN}{%
Costa Rica's National Artificial Intelligence Strategy (ENIA 2024-2027) establishes
principles of transparency and explainability, fairness and non-discrimination, and
accountability, but does not specify technical mechanisms to demonstrate compliance.
This article presents an operational framework that translates these three
principles into nine verifiable requirements, computational tests, and auditable
artifacts. The research followed a design-science approach: the requirements were
derived from the ENIA, the European AI Act, and the NIST AI Risk Management
Framework, and three verification modules were implemented: explainability through
SHAP; fairness through differences in true- and false-positive rates between
subgroups, with thresholds declared before measurement; and traceability through a
log that records each inference and links it to its explanation. The framework was
applied to a daily-activity classifier based on passive infrared (PIR) sensors,
using two distinct CASAS datasets: Aruba, with one home, and a series of nine homes
of older adults. In both pilots it produced a complete evidence dossier, with all
nine requirements covered, and the evaluated system received a non-compliance
verdict. Disparities appeared across time-of-day bands, day types and, above all,
homes, which account for 14 of the 29 combinations that exceeded the threshold in
the multi-home pilot. With ten seeds, the verdict held and 32 of the 33 combinations
above the threshold persisted in all of them. Results show that the selected
principles can be operationalized through reproducible tests, although the scarcity
of demographic metadata in public data limits the fairness evaluation.}

\newcommand{\palabrasclaveEN}{%
Artificial intelligence; algorithmic governance; AI regulatory compliance;
algorithmic fairness; explainability; passive infrared sensors; aging in place}
